%%%%%%%%%%%%%%%%%%%%%%%%%%%%%%%%%%%%%%%%%%%%%%%%%%%%%%%%%%%%
%% CHAPTER 4 -- From Letters to Vectors
%%%%%%%%%%%%%%%%%%%%%%%%%%%%%%%%%%%%%%%%%%%%%%%%%%%%%%%%%%%%

\chapter{From Letters to Vectors}
\label{ntg:ch4}

\section*{Setting the Scene}
\addcontentsline{toc}{section}{Setting the Scene}

\subsection*{The problem}

A computer cannot read words. It can read numbers. The neural
network we described in Chapter~3 takes a list of numbers as
input, performs arithmetic on them, and produces a list of
numbers as output. Before any of the apparatus of the previous
chapters can do anything useful with a piece of text, the text
has to be turned into numbers.

This chapter is about how that turning-into-numbers happens. It
proceeds in two steps. The first step is called
\emph{tokenization}: the text is chopped into pieces, called
\emph{tokens}, and each piece is given a number. The second
step is called \emph{embedding}: each token's number is used to
look up a list of further numbers, called the token's
\emph{vector} or \emph{embedding}, that captures something about
the token's meaning. The two steps sound prosaic. They are
also the place where several of the most user-visible quirks
of modern Large Language Models come from --- the way the model
miscounts letters in a word, the way the same sentence in two
languages costs different amounts to process, the way bias in
the training corpus survives into the model's behaviour.

\subsection*{How we got here}

The earliest computational treatments of language took the
\emph{word} as the unit of analysis. This was intuitive --- a
dictionary is a list of words --- and almost immediately
catastrophic. English has too many words. Inflectional
languages like Finnish and Hungarian have far too many.
Compound-forming languages like German and Dutch can produce a
new word any time you want one. Every model encountered, in
production, words that were not in its vocabulary. The standard
response was to replace the unknown word with a single special
symbol, conventionally written ``\texttt{<unk>}'', that
destroyed any information the word had carried. A model whose
answer to half the questions it received was effectively
``unknown'' was not a useful model.

For a brief period in the mid-2010s the field considered the
opposite extreme: take the \emph{character} as the unit. Every
text becomes a sequence of letters. There is no out-of-
vocabulary problem because every letter is in the vocabulary.
The trouble is that a single English word becomes five or six
units, an entire paragraph becomes thousands, and the network
spends most of its capacity learning to recognise that
``c'', ``a'', and ``t'' together mean ``cat'' --- a thing it
should be able to take for granted. Character-level models were
elegant. They were, on the whole, defeated by the practical
costs.

The compromise that won was the \emph{subword}: a unit smaller
than a word but larger than a letter, learned from the corpus
rather than handed down by a linguist. The algorithm in
predominant use today, \emph{Byte-Pair Encoding}, was originally
designed by Philip Gage in 1994 as a data compression technique.
Rico Sennrich, Barry Haddow, and Alexandra Birch, at the
University of Edinburgh in 2016, adapted it for neural machine
translation. Their idea was simple: start with the alphabet,
look through the corpus for the most common pair of adjacent
units, and merge that pair into a single new unit. Repeat. After
a few thousand merges, the resulting vocabulary contains
common words as single units, common parts of words as units,
and rare words decomposed into their parts. The vocabulary
size is fixed --- around thirty to one hundred thousand units
in modern models --- and any string in any language can be
represented in it without an out-of-vocabulary problem.

Two close cousins of byte-pair encoding share the field.
WordPiece was developed by Mike Schuster and Kaisuke Nakajima
at Google for voice search; it became famous when Google's BERT
model used it in 2018. SentencePiece, by Taku Kudo also at
Google, handles languages without obvious word boundaries
(Chinese, Japanese, Thai) cleanly. By 2020 every modern
language model used one of these three algorithms. The
differences between them are mostly engineering details. The
qualitative behaviour --- chop text into subword units, give
each unit a number, never have an out-of-vocabulary case --- is
the same.

The second step, the embedding, has its own history. The
foundational idea is older than computational linguistics:
\emph{the meaning of a word can be inferred from the company
it keeps}. The British linguist J.~R.~Firth said this in 1957.
The American linguist Zellig Harris, Chomsky's teacher, said
something close to it in 1954. The principle is that two
words that appear in similar contexts probably mean similar
things, and the corollary is that you can characterise the
meaning of a word by recording the contexts in which it
appears. The principle was used in vector-space models of
information retrieval as early as the 1960s.

The decisive computational step came in 2003 with Yoshua
Bengio's neural language model and in 2013 with Tom\'a\v{s}
Mikolov's \emph{word2vec}. Mikolov simplified Bengio's idea
into a system that could be trained on enormous corpora cheaply,
and the resulting word vectors --- lists of typically three
hundred numbers per word --- had a striking property. The
geometry of the vector space encoded meaning. Words with
similar meanings were close to one another in the space. The
direction from ``man'' to ``woman'' was approximately the same
as the direction from ``king'' to ``queen'', so that the
arithmetic ``king minus man plus woman'' landed near
``queen''. The geometry was not perfect --- the analogy trick
broke down on most analogies that were not specifically about
gender or capital cities --- but it was a real and reproducible
phenomenon, and it convinced the field that meaning could be
represented numerically.

Modern Large Language Models do not, in general, train word
vectors separately and then plug them into the network. They
train the embedding as the bottom layer of the same neural
network we described in Chapter~3, using the same training
procedure of Chapter~2. The geometric properties Mikolov
observed remain present, in modified form, in the embeddings of
modern models, and they continue to be one of the things that
makes the language modelling pipeline work.

\subsection*{Where we stand}

A modern tokenizer is a small, deterministic, frozen artefact.
It is trained once, before the main model training begins, on a
representative sample of the corpus. After that, it is fixed for
the life of the model: changing it would invalidate every token
the model had ever seen. The tokenizer is, in essence, a lookup
table from strings to integers and back. Given a sentence, it
produces a list of integers (the token IDs). Given the list of
integers, it reconstructs the sentence almost exactly --- modulo
a few normalisation choices about whitespace and Unicode form.

The embedding layer takes each token ID and looks up the
corresponding vector --- the list of numbers, typically a few
thousand long, that the trained model uses to represent that
token internally. The vectors live in a high-dimensional space
in which directions and distances encode meaning, in a way the
training procedure discovered without anyone telling it to. The
rest of the model --- the deep stack of layers from Chapter~3
plus the attention machinery of Chapter~5 --- operates on these
vectors and produces, at the end, a probability distribution
over which token should come next.

\subsection*{What goes wrong}

Tokenization is a quiet source of several user-visible quirks.

The most consequential is the multilingual disparity. The same
sentence in Thai, Amharic, or Burmese costs typically two to
five times as many tokens as the English original under a
tokenizer trained predominantly on English. This means a
context window of a fixed token length covers far less of those
languages than of English, and an API priced per token charges
far more for the same content. There is no clean fix at the
model level; partial fixes include training tokenizers on more
balanced corpora, which several frontier model providers have
begun to do.

A second issue is the ``model doesn't know X'' phenomenon. The
model often does know X. The token sequence that names X is
just rare in the training data, or it is split awkwardly across
token boundaries, and the model's grasp of it is correspondingly
weak. Calling this an ignorance failure obscures the underlying
tokenization issue. We will return to this in the engineering
section below.

A third issue is brittleness at token boundaries. Two prompts
that look identical to a human can segment differently if
whitespace, capitalisation, or Unicode normalisation differs by
even a single character. The model behaves differently on the
two segmentations, sometimes dramatically. Engineers who paste
prompts together at runtime, without testing the resulting
tokenization, ship products that misbehave in ways no test
caught.

A fourth issue is adversarial. Unicode is permissive about
characters that look identical to a human reader (``homoglyphs'',
like the Latin and Cyrillic letter ``a'') and about invisible
code points (zero-width spaces, bidirectional override
characters). An attacker can craft a prompt whose visible text
says one thing and whose tokenized form says another. We will
see this attack class properly in Chapter~12.

\section{The idea, in plain language}
\addcontentsline{toc}{section}{The idea, in plain language}

\subsection*{Tokenization, by hand}

The cleanest way to feel what tokenization does is to walk
through it on a single sentence.

Take the sentence ``The cat sat on the windowsill.'' A modern
tokenizer might segment this into the tokens \texttt{The},
\texttt{ cat}, \texttt{ sat}, \texttt{ on}, \texttt{ the},
\texttt{ window}, \texttt{sill}, \texttt{.}. Notice three
things. First, the leading space is part of the token: the
tokenizer treats ``\texttt{ cat}'' (with space) and
``\texttt{cat}'' (without) as different tokens, because they
behave differently. Second, ``windowsill'' has been split into
two tokens, ``\texttt{ window}'' and ``\texttt{sill}''. The
tokenizer's vocabulary contained ``\texttt{ window}'' as a
common word but not the full compound, so the compound was
decomposed. Third, the period is its own token. Punctuation
almost always gets its own tokens.

Now consider the same sentence in Thai. The exact Thai
characters do not matter for the example; the point is that a
sentence of similar meaning in Thai, Burmese, or Amharic might
be segmented into fifteen or twenty tokens rather than eight,
simply because the tokenizer was trained predominantly on
English text and never saw enough Thai to recognise common
Thai sequences as tokens worth giving their own slot. The
model can still process the Thai sentence, but it spends more
of its budget doing so, and some of the relationships the
model would learn easily in English
(``\texttt{ window}'' as a single semantic unit) are
fragmented across many small tokens it has to reassemble.

The same fragmentation explains why models cannot reliably count
the letters in a word. The word ``strawberry'' may be tokenized
as ``\texttt{straw}'' and ``\texttt{berry}'', or as a single
token, depending on the tokenizer. The model never sees the
individual letters. Asking it ``how many R's are in
strawberry'' is asking it about a property of the spelled-out
form that is not part of its input. Sometimes it gets the
answer right because something in its training data discussed
the spelling. Often it does not. This is not a deep failing of
language modelling. It is a property of the tokenization step.

\subsection*{Embeddings as a map of meaning}

The embedding step takes each token ID and looks up the list of
numbers (typically two to four thousand of them in modern
models) that the model uses to represent that token internally.
You can think of each embedding as a coordinate in a very
high-dimensional space. We cannot picture a four-thousand-
dimensional space, but we can carry several useful intuitions
from lower-dimensional analogies.

In two dimensions: imagine plotting the words of English on a
sheet of paper, with related words near one another. Words for
animals would form a region. Within it, ``dog'' would be near
``cat'' and far from ``elephant''. The arrangement is what
linguists call a \emph{semantic field}. The field is not a
discrete category; it is a continuous neighbourhood, and a word
can be in several overlapping neighbourhoods at once.

In four thousand dimensions: there are enough directions that
each different ``meaning'' a word has can be encoded along a
different one. The direction from ``man'' to ``woman'' is one
direction. The direction from ``Paris'' to ``France'' is
another. The direction that distinguishes ``ran'' from ``run''
is a third. The training procedure of Chapter~2 worked these
directions out from the contexts in which the words appeared,
without being told what any of them meant.

The magic of word2vec --- ``king minus man plus woman is near
queen'' --- is the geometric statement that these directions
exist as roughly straight lines in the embedding space and that
they roughly add. This is not exactly true. The directions are
not exactly straight, the analogies are not exactly preserved,
and a long literature has grown up showing where the trick
fails. But it is true \emph{enough} to be useful, and the
overall fact that words with similar meaning end up near each
other in the embedding space is what makes the rest of the
model's machinery work. The model does not need a separate
look-up table that says ``cat is similar to kitten''. It can
read the similarity off the geometry.

\subsection*{Why training pressure carves meaning into the space}

The claim that words appearing in similar contexts end up with
similar vectors is not a design decision. It is a consequence.
The training procedure makes it happen without being told to, and
understanding why is the difference between taking the geometry
on faith and actually seeing how it arises.

Consider two words: ``dog'' and ``cat''. Both appear, in English
text, after sentences like ``I have a'', ``she is allergic to
her neighbour's'', and ``the vet said the''. Both appear before
words like ``food'', ``collar'', ``chased the''. From a next-token
prediction standpoint, these two words are nearly
interchangeable: knowing that the next word is ``dog'' tells you
roughly the same things about what will follow as knowing it is
``cat''.

Now suppose the model has assigned ``dog'' and ``cat'' very
different vectors. To predict well after ``dog'', it learns one
set of patterns. To predict well after ``cat'', it learns a
separate set of the same patterns. The same work is done twice.
If instead the model assigns them similar vectors, both words
arrive at the attention and feed-forward layers looking roughly
alike, and the model only has to learn the relevant patterns
once to predict correctly in both cases. The training
pressure --- the relentless push to reduce prediction error
across billions of examples --- rewards the second arrangement.
The gradient updates accumulate in the direction of similar
vectors for similarly-behaving words. Nobody specified this.
It is what minimising the loss looks like.

This is the distributional hypothesis made mechanistic. Harris
and Firth observed the correlation: similar contexts, similar
meaning. The training procedure explains the mechanism: similar
contexts create shared prediction targets, and the cheapest way
to meet those targets is to produce shared internal
representations. Over enough training examples, that
pressure is large enough to carve the geometry the word2vec
demonstration made famous. The direction from ``man'' to
``woman'', the direction from ``Paris'' to ``France'', the
cluster of animal words near one another --- none of these were
specified. All of them were the path of least resistance for a
system trained to predict text.

One consequence worth noting: the geometry encodes whatever
statistical regularities were present in the corpus, whether
or not those regularities reflect the world accurately. A corpus
that consistently associates ``nurse'' with female pronouns will
produce a geometry in which ``nurse'' points in the direction
of ``she'' more than ``he''. That is not the model being biased
in some separate sense; it is the model faithfully encoding what
the corpus said. The geometry is a mirror, not a window.

\subsection*{What the embedding is and is not}

It is worth being careful about what the embedding does and
does not contain.

It contains, encoded in the geometry, the distributional
statistics of the training corpus. Words that appeared in
similar contexts in the corpus end up near one another in the
space. Directions in the space sometimes correspond to
interpretable axes of meaning (gender, tense, formality), often
do not, and almost always vary in interpretability across
models.

It does not contain anything about the world that the corpus
did not contain. If the corpus did not discuss a particular
chemical compound, the model has no embedding for it that
encodes anything useful about its properties. (The model may
still have some embedding for it, if the chemical name was in
the corpus at all. The embedding will, however, be sparse,
and the model's behaviour around it correspondingly fragile.)

It does not contain truth values. The embedding for ``Earth''
and the embedding for ``flat'' have a relationship in the space
that reflects how often they appeared together in the corpus,
not whether the Earth is in fact flat. Models trained on
corpora that contain a lot of conspiracy theory will have
embeddings that reflect the conspiracy theory's statistical
prominence, not its truth. This is one of the reasons curating
the training corpus matters.

\section{The engineering consequence}
\addcontentsline{toc}{section}{The engineering consequence}

\paragraph{Token counts are the unit of cost.} Almost every
modern LLM API charges by the token, both for input (the prompt
you send) and for output (the response you receive). A user
who sends ten thousand words of context is paying for the
tokenization of those ten thousand words. The same content in
different languages costs different amounts. The same content
expressed concisely versus verbosely costs different amounts.
For a product team running an LLM at scale, the tokenization
step has a direct line to the monthly bill.

\paragraph{The context window is in tokens.} Modern models have
a maximum context length, expressed in tokens, beyond which they
either truncate the input or refuse the request. In 2024 this
was typically between 8{,}000 and 200{,}000 tokens. The number
of words this corresponds to depends on the language, the
tokenizer, the proportion of code, and the proportion of
unusual content. A team designing a product around an LLM has to
budget for the tokenization rather than for the visible text.
When the budget is tight, this is one of the first places to
optimise.

\paragraph{The tokenizer is a frozen artefact.} For any given
model checkpoint, the tokenizer cannot be changed without
re-training the model. This means that if a tokenizer was
trained predominantly on English in 2020, the model still
suffers from that imbalance in 2024. Newer models use newer
tokenizers, and the gap between an old and a new model can
include differences in tokenization that the model cards do not
always foreground.

\paragraph{Embedding similarity is a useful primitive.} Because
the embedding space encodes meaning geometrically, the
similarity between two pieces of text can be measured by
embedding both and computing how close they are in the space.
This is the basis of the \emph{vector database} that
underpins the retrieval systems of Chapter~10. It is not a
perfect measure of meaning --- two pieces of text can be close
in the embedding space without being equivalent in any deep
sense --- but it is a much better measure than the keyword
matching it replaced.

\section{What goes wrong, and why}
\addcontentsline{toc}{section}{What goes wrong, and why}

\paragraph{Letter-level operations.} The model does not see
letters; it sees tokens. Asking it to count the letters in a
word, to spell a word backwards, to pick a word that starts
with a particular letter, or to perform any operation that
requires inspecting individual characters within a token will
sometimes succeed (because the training corpus discussed the
spelling) and will sometimes fail in surprising ways. The
failure mode is not a failure of intelligence; it is a failure
of input representation. Engineers building products that
depend on letter-level operations should not use the bare
language model for them.

\paragraph{Multilingual cost and quality.} A model trained
predominantly on English text is, on average, less fluent in
non-English languages. The disparity is partly in the model
weights (more training compute was spent on English) and partly
in the tokenization (the tokenizer represents English more
efficiently than most other languages). The combined effect is
that the same task costs more and is performed worse in less
well-represented languages. This is a known and unevenly
addressed inequity, and a non-technical reader evaluating an
LLM product for a non-English deployment should ask the vendor
for performance numbers in the target language specifically.

\paragraph{Bias as inherited corpus statistics.} Words associated
in the training corpus with a particular gender, ethnicity, or
profession will, in the model's embeddings, encode that
association. A model whose corpus contained ``the doctor said
he'' more often than ``the doctor said she'' will, on average,
predict masculine pronouns after ``the doctor said''. This is
not a separate failure of the model. It is a faithful encoding
of the statistics of the corpus. Reducing the bias is possible
and is the subject of an ongoing literature, but doing so
without compromising other aspects of model behaviour is hard,
and the trade-offs are usually not transparent in vendor
documentation.

\paragraph{The hidden interface between text and tokens.} The
gap between what a user sees on the screen and what the model
sees in its input is invisible to most users and to many
engineers. This is the surface on which several attacks
operate. We will see in Chapter~12 the class of attacks that
exploit invisible Unicode characters, look-alike letters, and
unexpected normalisations to make the model see one thing while
the user sees another. The tokenization step is the place where
the gap is opened.

\section*{Bridges}
\addcontentsline{toc}{section}{Bridges}

This chapter built the bridge from text to numbers that
everything else in the book takes for granted. Chapter~5 takes
the embedded tokens as the input to attention. Chapter~6
trains the embedding layer along with the rest of the model and
gives us the larger context of pretraining. Chapter~10 uses the
embedding space directly as the basis for retrieval. Chapter~12
returns to the tokenization gap as the surface on which several
classes of attack operate. The next chapter --- \emph{Attention,
or How a Machine Reads} --- shows what the model does with the
embedded tokens once it has them.
